\section{Conclusion}


Post-quantum migration is a process, not an event. The three-phase strategy---hybrid, PQ-primary, pure PQ---allows a production blockchain network to progressively strengthen its cryptographic foundations without coordinated flag-day upgrades. The Lux network's Phase 1 deployment (triple-proof consensus with BLS + ML-DSA + Pulsar) provides hybrid security from genesis. Phases 2 and 3 are defined, with concrete trigger conditions tied to NIST standardization, government mandates, and quantum computing developments. At no point during the migration is the network less secure than it was before---each phase strictly increases the post-quantum security level while maintaining backward compatibility with previous-phase clients.

\begin{thebibliography}{9}
\bibitem{fips203} NIST, ``FIPS 203: Module-Lattice-Based Key-Encapsulation Mechanism Standard (ML-KEM),'' August 2024.
\bibitem{fips204} NIST, ``FIPS 204: Module-Lattice-Based Digital Signature Standard (ML-DSA),'' August 2024.
\bibitem{fips205} NIST, ``FIPS 205: Stateless Hash-Based Digital Signature Standard (SLH-DSA),'' August 2024.
\bibitem{cnsa2} NSA, ``CNSA 2.0: Cybersecurity Advisory on Post-Quantum Cryptography,'' 2022.
\bibitem{ir8547} NIST, ``IR 8547: Transition to Post-Quantum Cryptography Standards,'' 2024.
\bibitem{lux-triple} Lux Industries, Inc., ``Triple-Proof Quantum Finality: Post-Quantum Consensus for the Lux Network,'' 2026.
\bibitem{lux-pq} Lux Industries, Inc., ``Post-Quantum Cryptographic Suite for EVM,'' 2024.
\bibitem{lux-corona} Lux Industries, Inc., ``Corona: Lattice-Based Threshold Ring Signatures for Anonymous Consensus,'' 2025.
\end{thebibliography}

